\scanpage{26}
sous l'hypothèse anabélienne faite, et celle de dominance. (Cette deuxième n'est pas superflue, comme on voit en prenant \deleted{\unsure{pour} $Y$ \inserted{$=G\backslash Z$, où $G$ est} un groupe elliptique et $Z$ un sous-groupe fini $\unsure{+1|\ne 1}$, et pour $f,g$} \inserted{deux applications constantes de $X$ dans $Y$!})

Tout d'abord, on vérifie qu'une variété élémentaire d'Artin se plonge bien \inserted{par un $i:Y\longrightarrow G$} dans une ext. d'une V.A. par un tore -- donc on aura $(ig)=\tau_u(i\circ f)$. \deleted{Donc \unsure{tout}} \inserted{D'autre part,} \deleted{Quitte à remplacer} $f$,\deleted{$g$} étant dominante, il existe une partie ouverte \inserted{non vide} $\unsure{V|W}\subset Y$ contenue dans $\operatorname{Im}f$, et on aura donc $\unsure{V|W}+u\subset Y$. Soit $U\subset X$ un ouvert non vide tel que $f(U)\subset V$ \deleted{$g(U)\subset V$} \deleted{\unsure{donc qui} $\unsure{Q}$} quitte à remplacer $X$ par $U$, on est ramené au cas où $f$\deleted{,$g$} envoie $X$ dans un ouvert $V$ de $Y$ tel que $V+u\subset Y$, et $g\circ i=\tau_u(i\circ f)$. On \deleted{\unsure{sait}} \inserted{voit} que $\pi_1(f)$ est un composé (où $f':X\unsure{-|\longrightarrow}U$ est \deleted{$\unsure{f}$} induit par $f$ et $j:U\unsure{-|\longrightarrow}Y$ est l'inclusion)
% Each empty editorial doubt marker below a dashed shaft concerns only its terminal head: dash first, right arrow alternative. Objects and source labels are not thereby made doubtful.
\diagram{\begin{tikzcd}[ampersand replacement=\&]\pi_1(X) \arrow[r,dash,"\pi_1(f')","{\unsure{}}"'] \& \pi_1(U) \arrow[r,"\pi_1(j)"] \& \pi_1(Y)\end{tikzcd}}
tandis que $\pi_1(g)$ est le composé
\diagram{\begin{tikzcd}[ampersand replacement=\&]\pi_1(X) \arrow[r,dash,"\pi_1(f')","{\unsure{}}"'] \& \pi_1(U) \arrow[r,dash,"{\pi_1(\unsure{j'|g'})}","{\unsure{}}"'] \& \pi_1(Y)\end{tikzcd}}
% The overwritten j'/g' label is one unresolved glyph; the visible overtracing does not securely distinguish a cancelled g from a replacement j.
où $\deleted{\unsure{j}}$ $j':U\unsure{-|\longrightarrow}Y$ est induit par $\tau_u$.
\margin{25}
